\documentclass[12pt,oneside]{book}

%% ---------------------------------------------------------------
%% Оригинал-макет по техническим требованиям ИПЦ НГУ:
%% формат А5 (60x84 1/16), 148x210 мм; полоса набора 115x175 мм;
%% поля: верхнее, правое, левое 17 мм, нижнее 20 мм;
%% Times New Roman 10-11 пт, одинарный интервал, абзац 5 мм,
%% колонцифра внизу по центру, автоматические переносы,
%% выравнивание по ширине.
%% ---------------------------------------------------------------

%% А4 (60x84 1/8): страница 210x297 мм, полоса набора 170x257 мм,
%% поля по 20 мм. Нижнее поле берём 20 мм плюс колонцифра в footskip.
\usepackage[paperwidth=210mm,paperheight=297mm,
            top=20mm,left=20mm,right=20mm,bottom=20mm,
            footskip=12mm,headsep=6mm]{geometry}

\usepackage{fontspec}
\usepackage{polyglossia}
\setmainlanguage{russian}
\setotherlanguage{english}

\setmainfont{Times New Roman}[Ligatures=TeX]
\setsansfont{Arial}
\setmonofont{Menlo}[Scale=0.78]
%% polyglossia требует явно указать шрифты для кириллицы
\newfontfamily\cyrillicfont{Times New Roman}[Ligatures=TeX]
\newfontfamily\cyrillicfontsf{Arial}
\newfontfamily\cyrillicfonttt{Menlo}[Scale=0.78]

\usepackage{microtype}
\usepackage{unicode-math}
% математический шрифт, парный к Times New Roman
\setmathfont{texgyretermes-math.otf}

\usepackage{graphicx}
%% таблицы (набор макросов, который ожидает pandoc)
\usepackage{longtable,booktabs,array}
\newcounter{none}                  % для ненумерованных таблиц
\usepackage{calc}                  % расчёт ширины колонок
\usepackage{etoolbox}
\makeatletter
\patchcmd\longtable{\par}{\if@noskipsec\mbox{}\fi\par}{}{}
\makeatother
\IfFileExists{footnotehyper.sty}{\usepackage{footnotehyper}}{\usepackage{footnote}}
\makesavenoteenv{longtable}
%% таблицы данных набираем мельче — в полосе 115 мм иначе не помещаются
\AtBeginEnvironment{longtable}{\footnotesize}
\usepackage{caption}
\usepackage[table]{xcolor}
\usepackage{fancyvrb}
\usepackage{fvextra}   % перенос длинных строк в листингах
\usepackage{framed}
\usepackage{titlesec}
\usepackage{fancyhdr}
\usepackage{enumitem}
\usepackage{hyperref}

%% --- предметный указатель ---------------------------------------
%% Сортировку кириллицы делает xindy (makeindex её не умеет),
%% вызов — в build.sh между проходами xelatex.
\usepackage[noautomatic]{imakeidx}
\makeindex[title={Предметный указатель}, columns=2, intoc]

%% --- набор -------------------------------------------------------
\linespread{1.0}
\setlength{\parindent}{10mm}
\setlength{\parskip}{0pt plus 1pt}
\sloppy
\hyphenpenalty=50
\tolerance=2000
\emergencystretch=3em
\clubpenalty=10000
\widowpenalty=10000

%% --- колонтитулы: колонцифра внизу по центру ---------------------
\pagestyle{fancy}
\fancyhf{}
\fancyfoot[C]{\small\thepage}
\renewcommand{\headrulewidth}{0pt}
\renewcommand{\footrulewidth}{0pt}
\fancypagestyle{plain}{%
  \fancyhf{}\fancyfoot[C]{\small\thepage}%
  \renewcommand{\headrulewidth}{0pt}}

%% --- заголовки ---------------------------------------------------
\titleformat{\chapter}[display]
  {\normalfont\Large\bfseries\raggedright}
  {\normalsize\MakeUppercase{\chaptertitlename}\ \thechapter}{6pt}{\Large}
\titlespacing*{\chapter}{0pt}{0pt}{10pt}

\titleformat{\section}{\normalfont\large\bfseries\raggedright}{\thesection}{0.6em}{}
\titlespacing*{\section}{0pt}{9pt plus 2pt}{4pt}
\titleformat{\subsection}{\normalfont\normalsize\bfseries\raggedright}{\thesubsection}{0.6em}{}
\titlespacing*{\subsection}{0pt}{7pt plus 2pt}{3pt}
\titleformat{\subsubsection}{\normalfont\normalsize\itshape\raggedright}{}{0em}{}
\titlespacing*{\subsubsection}{0pt}{6pt plus 1pt}{2pt}

\setcounter{secnumdepth}{3}
\setcounter{tocdepth}{2}

%% --- листинги ----------------------------------------------------
%% макросы подсветки синтаксиса, которые генерирует pandoc
$highlighting-macros$

\definecolor{shadecolor}{gray}{0.955}
%% код набираем мелко и с переносом длинных строк — полоса набора узкая
\DefineVerbatimEnvironment{Highlighting}{Verbatim}%
  {commandchars=\\\{\},fontsize=\footnotesize,xleftmargin=3mm,
   breaklines,breakanywhere,breaksymbolleft={}}
\DefineVerbatimEnvironment{verbatim}{Verbatim}%
  {fontsize=\footnotesize,xleftmargin=3mm,breaklines,breakanywhere,
   breaksymbolleft={}}
\renewenvironment{Shaded}{\begin{snugshade}}{\end{snugshade}}
\setlength{\FrameSep}{1.5pt}
\setlength{\OuterFrameSep}{1.5pt}

%% --- прочее ------------------------------------------------------
\setlist{nosep,leftmargin=5mm,topsep=2pt,partopsep=0pt}
\captionsetup{font=small,skip=4pt}
%% Картинки вписываем в полосу набора: не шире 92% строки
%% и не выше 42% полосы, чтобы иллюстрация не занимала страницу целиком.
\setkeys{Gin}{width=0.92\linewidth,height=0.30\textheight,keepaspectratio}

%% \pandocbounded — макрос pandoc 3.x, масштабирует картинку под полосу
\makeatletter
\newsavebox\pandoc@box
\newcommand*\pandocbounded[1]{%
  \sbox\pandoc@box{#1}%
  \Gscale@div\@tempa{0.30\textheight}{\dimexpr\ht\pandoc@box+\dp\pandoc@box\relax}%
  \Gscale@div\@tempb{0.92\linewidth}{\wd\pandoc@box}%
  \ifdim\@tempb\p@<\@tempa\p@\let\@tempa\@tempb\fi%
  \ifdim\@tempa\p@<\p@\scalebox{\@tempa}{\usebox\pandoc@box}%
  \else\usebox{\pandoc@box}%
  \fi%
}
\makeatother

\hypersetup{
  colorlinks=false, pdfborder={0 0 0},
  pdftitle={$title$}, pdfauthor={$author$},
  bookmarksnumbered=true, bookmarksopen=true}
\urlstyle{same}

% длинные url и пути не должны вылезать за полосу набора
\usepackage{xurl}

\providecommand{\tightlist}{\setlength{\itemsep}{0pt}\setlength{\parskip}{0pt}}

%% Команды, которые понимает MathJax на сайте, но не знает LaTeX
\providecommand{\lt}{<}
\providecommand{\gt}{>}

%% Пустой квадратик для чек-листов (- [ ] в markdown).
%% amssymb здесь не подключить — он конфликтует с unicode-math.
\providecommand{\square}{}
\renewcommand{\square}{\mbox{\raisebox{0.1ex}{\framebox[1.5ex]{\rule{0pt}{1.1ex}}}}}

\begin{document}

%% ================= ТИТУЛЬНЫЙ ЛИСТ =================
\thispagestyle{empty}
\begin{center}
  \setlength{\parindent}{0pt}
  {\small МИНИСТЕРСТВО НАУКИ И ВЫСШЕГО ОБРАЗОВАНИЯ РФ}\\[2pt]
  {\small НОВОСИБИРСКИЙ ГОСУДАРСТВЕННЫЙ УНИВЕРСИТЕТ}\\[2pt]
  {\small Физический факультет}
  \vfill
  {\large В.\,В.~Федоров}\\[16mm]
  {\LARGE\bfseries РАЗРАБОТКА И ПРИМЕНЕНИЕ}\\[4pt]
  {\LARGE\bfseries ПРОГРАММНОГО ОБЕСПЕЧЕНИЯ}\\[4pt]
  {\LARGE\bfseries В ФИЗИЧЕСКИХ ИССЛЕДОВАНИЯХ}\\[8mm]
$if(volume)$
  {\Large\bfseries Часть $volume$}\\[3pt]
  {\large\itshape $volumetitle$}\\[8mm]
$endif$
  {\large Учебное пособие}
  \vfill
  {\small Новосибирск}\\
  {\small 2026}
\end{center}
\clearpage

%% ================= ОБОРОТ ТИТУЛА =================
\thispagestyle{empty}
\setlength{\parindent}{0pt}
{\small
УДК~004.4:53(075.8)\\
ББК~32.973-018я73\\
Ф333
\vspace{6mm}

% Рецензент — член-корреспондент РАН. Формат записи взят из образца
% выходных данных ИПЦ: «член-корр. РАН, проф. В. Г. Романов».
% Вторая строка оставлена на случай, если ИПЦ потребует двух рецензентов:
% образец выходных данных даёт «Рецензент» в единственном числе, а
% выписка из протокола кафедры говорит о «рецензиях» во множественном.
% Если хватит одного — убрать \textit{Рецензенты} на \textit{Рецензент}
% и снять вторую строку.
\textit{Рецензенты}\\
чл.-корр. РАН \underline{\hspace{42mm}}\\[1pt]
\underline{\hspace{55mm}}\\[1pt]
{\footnotesize (учёная степень, звание, И.\,О.~Фамилия --- заполнить)}
\vspace{4mm}

Издание подготовлено в рамках реализации Программы развития
Новосибирского государственного университета.
\vspace{4mm}

\textbf{Федоров, В.\,В.}\\
Ф333\quad Разработка и применение программного обеспечения в физических
исследованиях : учеб. пособие$if(volume)$ : в 2 ч. / В.\,В.~Федоров ;
Новосиб. гос. ун-т. --- Новосибирск : ИПЦ НГУ, 2026. ---
Ч.~$volume$: $volumetitle$. --- \pageref{LastPage}~с.$else$ /
В.\,В.~Федоров ; Новосиб. гос. ун-т. --- Новосибирск : ИПЦ НГУ, 2026.
--- \pageref{LastPage}~с.$endif$
\vspace{3mm}

ISBN 978-5-4437-XXXX-X$if(volume)$ (ч.~$volume$)\\
ISBN 978-5-4437-YYYY-Y$endif$
\vspace{4mm}

$abstract$

Предназначено для студентов второго курса, обучающихся по направлению
<<Прикладная математика и физика>>, в рамках дисциплины <<Применение
Python в физических исследованиях>>, и будет полезно всем, кто
разрабатывает программное обеспечение для научных исследований.
\vspace{4mm}

УДК~004.4:53(075.8)\\
ББК~32.973-018я73
\vspace{5mm}

Рекомендовано к печати кафедрой \underline{\hspace{45mm}}\\
(протокол №~\underline{\hspace{12mm}} от \underline{\hspace{20mm}})
\vfill

\noindent
\begin{tabular}{@{}l@{\hspace{6mm}}l@{}}
                       & \copyright{} Новосибирский государственный\\
ISBN 978-5-4437-XXXX-X$if(volume)$ (ч.~$volume$)$endif$ & \phantom{\copyright{}} университет, 2026\\
                       & \copyright{} В.\,В.~Федоров, 2026\\
\end{tabular}
}
\clearpage
\setlength{\parindent}{10mm}

%% ================= ОГЛАВЛЕНИЕ =================
\pagestyle{fancy}
\setcounter{page}{3}
\renewcommand{\contentsname}{Оглавление}
\tableofcontents
\clearpage

%% ================= ТЕКСТ =================
$body$

\clearpage
\printindex

\clearpage
\thispagestyle{empty}
\vspace*{\fill}
\begin{center}
{\small
Учебное издание
\vspace{6mm}

\textbf{Федоров Вячеслав Васильевич}
\vspace{4mm}

РАЗРАБОТКА И ПРИМЕНЕНИЕ ПРОГРАММНОГО ОБЕСПЕЧЕНИЯ\\
В ФИЗИЧЕСКИХ ИССЛЕДОВАНИЯХ
$if(volume)$

Часть $volume$. $volumetitle$
$endif$
\vspace{4mm}

Учебное пособие
\vspace{10mm}

Редактор \underline{\hspace{30mm}}\\
Оригинал-макет \underline{\hspace{30mm}}\\
Обложка \underline{\hspace{30mm}}
\vspace{8mm}

\begin{tabular}{@{}l@{}}
Подписано в печать \underline{\hspace{20mm}}\quad Формат 60\,$$\times$$\,84 1/8\\
Уч.-изд. л. \underline{\hspace{12mm}}\quad
Усл. печ. л. \underline{\hspace{12mm}}\\
Тираж \underline{\hspace{12mm}} экз. \quad Заказ № \underline{\hspace{12mm}}\\
\end{tabular}
\vspace{8mm}

Издательско-полиграфический центр НГУ\\
630090, Новосибирск, ул. Пирогова, 2\\
ed@post.nsu.ru\\
https://www.nsu.ru/n/university/publishing-center/
}
\end{center}
\label{LastPage}
\end{document}
